%% kapa: exam=1 round=19 year=2008 period=2 form=A subject=english subject_ko=영어 seq=12 number=52
%% source: Q-Net 감정평가사 기출문제 (cst00309) artlSeq=1001962 file=extracted/1008/1001962/2008년 제19회 감정평가사 1차시험 기출문제/2교시.hwp
%% conversion: kordoc 4.24.51 -> tools/gen1.py (auto); review=proofread-sonnet; flags=none
\begin{question}{52}
% [공통 지문/지시문: 원문에서 여러 문항에 걸쳐 제시됨]
\begin{framedtext}
※ 보기의 밑줄 친 표현과 가장 가까운 의미로 쓰인 것을 고르시오. (51-53) \par
\end{framedtext}
He chiseled the dragon's nostrils to flare and created fiery-looking eyes. He used pneumatic tools to carve the hooked beak, \uline{scales}, and fang-like teeth for maximum control.
\begin{choices}
  \item He compared the different \uline{scales} of the two tests.
  \item The map, on a \uline{scale} of 1:10,000, shows over 5,000 individual paths.
  \item Do not remove the \uline{scales} from the skin of the fish.
  \item The breakdown of law and order could result in killing on a massive \uline{scale}.
  \item The mechanic was ordered to fix twenty weighing \uline{scales} a day.
\end{choices}
\end{question}
\answer{3}
